% paper/sections/08_collocate.tex
\section{圧力・速度連成：Balanced--Force CCD 理論}
\label{sec:collocate}%
\label{sec:pressure}% external references

\noindent\textbf{本章の位置付け：}
界面適合非一様格子（§~\ref{sec:grid}章で詳述）上に，
すべての流体変数（$\bu, p, \psi$）を同一格子点に配置する
\textbf{コロケート配置}を採用する．
本章の問いは二つである．第一に，変密度二相流で中間速度を
どの楕円問題により発散ゼロ空間へ射影するか．第二に，その射影で使う
圧力による面加速度と毛管ジャンプ補正量を，同じ離散空間で
釣り合わせるには何を揃えなければならないか．
以下では，連続射影の導出から始め，コロケート格子で生じる離散的な破綻を同定し，
最後に圧力ジャンプ標準経路を支える Balanced--Force 条件と
FCCD 面フラックス構成を与える．

\noindent\textbf{本章の構成：}
\begin{enumerate}
  \item \textbf{§~\ref{sec:collocate_layout}（コロケート格子の利点と課題）：}
        コロケート格子の利点と，$2\Delta x$ チェッカーボード問題の根本原因
        （圧力・速度デカップリング）を示す．解消策は §~\ref{sec:checkerboard_solution} で論じる．
  \item \textbf{§~\ref{sec:helmholtz}（射影法の基礎：等密度場）：}
        Helmholtz 分解による発散ゼロ射影の幾何学的意味を示す（等密度場での定式化）．
  \item \textbf{§~\ref{sec:projection_derivation}（変密度射影法の段階的導出）：}
        等密度 Helmholtz 分解を変密度場へ拡張し，
        可変係数 PPE
        $\bnabla\cdot(1/\rho\,\bnabla\delta p)
        = (1/\Delta t_{\mathrm{proj}})\bnabla\cdot\bu^*$ を導出する．
        さらに，PPE の発散だけでは圧力反力の面代表が決まらないことを示し，
        $G_A=-M_A^{-1}D_f^TW_p$ と $L_A=D_fG_A$ の随伴ペアを
        標準補正の出発点にする．
  \item \textbf{§~\ref{sec:checkerboard_solution}（解決策の概観）：}
        Balanced--Force 演算子整合と FCCD 面フラックス部分系による
        デカップリング解消の役割分担を整理する．
  \item \textbf{§~\ref{sec:balanced_force}（Balanced--Force 条件：同一評価位置・随伴整合原理）：}
        圧力による面加速度と毛管ジャンプ補正量を同じ評価位置・同じ随伴ペアで
        構成する設計思想を示す．CSF 体積力はこの原理を一括一流体で読む対照形であり，
        標準経路では圧力ジャンプ補正量として扱う．
        演算子不一致がもたらす寄生流れの定量的解析を与える．
  \item \textbf{§~\ref{sec:bf_failure_modes}（Balanced--Force 失敗モード）：}
        位置，随伴性，係数平均，界面ジャンプのどこが崩れると
        連続平衡が離散平衡として残らないかを F-1--F-5 として切り分ける．
  \item \textbf{§~\ref{sec:bf_seven_principles}（Balanced--Force 七原則）：}
        F-1--F-5 を防ぐための P-1--P-7 を定式化し，
        主解消原則と次数改善・予防原則の役割を分離する．
  \item \textbf{§~\ref{sec:fccd_bf_sub_system}（FCCD 面フラックス部分系）：}
        FCCD の面フラックス表現を用いて PPE，速度補正，HFE/GFM ジャンプ行，
        表面エネルギー仮想仕事に基づく毛管閉包と成分体積制約付き鞍点系を同じ面の評価範囲へ閉じる．
\end{enumerate}

\paragraph{本章の数式の読み方．}
射影法の式は，圧力を後から足す補正ではなく，
中間速度から「非圧縮速度として許されない成分」を抜き取る操作として読むとよい．
等密度の Helmholtz 分解では，$\bu^*$ を発散ゼロ成分と勾配成分に分け，
勾配成分を圧力反力として取り除く．
変密度・二相流ではこの直感に，面係数，界面ジャンプ，毛管仕事が加わる．
Balanced--Force 条件は，圧力の面加速度と毛管の面補正量が
別々の格子位置で釣り合ったふりをしないための約束である．
したがって本章の式では，圧力値そのものよりも，
発散 $D_f$，面勾配 $G_A$，面係数 $A_f$ が同じ空間で対になっているかを追う．

\subsection{コロケート格子の利点と課題}
\label{sec:collocate_layout}
本研究では，すべての物理量（速度 $\bu$, 圧力 $p$, レベルセット $\phi, \psi$）を同一の格子点（セル中心）に配置する\textbf{コロケート格子}（collocated grid）を採用する．この配置は，複雑な形状の境界条件を同一の離散空間で扱いやすく，高次精度スキーム（例えば本稿の CCD 法）をすべての変数に統一的に適用できるという利点を持つ．

しかし，コロケート格子には特有の課題が存在する．それは，\textbf{圧力と速度の数値的な分離}に起因する，非物理的な振動解（チェッカーボードパターン）が発生しうることである．

\paragraph{なぜチェッカーボード問題が起きるのか：1 次元の例}

1次元の運動量方程式（$x$方向）を考えよう．セル $i$ における圧力勾配を単純な中心差分で近似すると，
\[ \left(\pd{p}{x}\right)_i \approx \frac{p_{i+1} - p_{i-1}}{2\Delta x} \]
セル $i$ の運動量は隣接するセル $i+1$ と $i-1$ の圧力しか参照せず，$p_i$ 自身は寄与しない．
圧力場が $\ldots +1, -1, +1, -1, \ldots$ のような波長 $2\Delta x$ の振動を持つ場合，
どのセルでも $\frac{p_{i+1}-p_{i-1}}{2\Delta x} = 0$ と評価され，
\textbf{速度場は高周波の圧力振動を「感知」できない}．
その結果，物理的に意味のないギザギザの圧力解が誤差として残り続ける——
これが圧力速度デカップリング問題の本質である．

\paragraph{CCD を用いた場合のデカップリング}
中心差分ではなく CCD を用いる場合も，本質的に同じデカップリング問題が起きる．
CCD の周波数応答は $kh \to \pi$ のとき $|\hat{D}_k| \to 0$ であり，
$2h$ 波長成分が不可視であることは中心差分と本質的に同じである．
対策は §~\ref{sec:checkerboard_solution} で述べる．

% ═══════════════════════════════════════════════════════════════════════════════
\subsection{射影法の数学的基礎：Helmholtz 分解}
\label{sec:helmholtz}
% ═══════════════════════════════════════════════════════════════════════════════

射影法（Projection 法）~\cite{Chorin1968,Guermond2006}は，Helmholtz--Hodge 分解に基づく速度場の射影操作である．
本節ではまず等密度場で原理を示し，変密度への拡張は §~\ref{sec:projection_derivation} で行う．

\subsubsection{Helmholtz 分解と射影法}

任意の十分滑らかなベクトル場 $\bm{w}$ は，スカラーポテンシャル $q$ を用いて
\begin{equation}
  \bm{w} = \underbrace{\bu}_{\bnabla\cdot\bu=0}
           + \bnabla q
  \label{eq:helmholtz}
\end{equation}
と発散ゼロ成分と勾配場（回転なし成分）に一意分解できる（適切な境界条件の下で）．
Navier--Stokes 方程式を時間離散化して得られる中間速度 $\bu^*$（$\bnabla\cdot\bu^* \neq 0$）に
この分解を適用する：
\begin{equation}
  \bu^* = \bu^{n+1} + \bnabla q,
  \qquad \bnabla\cdot\bu^{n+1}=0
  \label{eq:projection_decomp}
\end{equation}
ここで $q \propto p^{n+1}$ が圧力として同定される．
すなわち射影法とは，$\bu^*$ を発散ゼロ速度空間（Leray 空間）へ直交射影する操作である．

\subsubsection{PPE の導出}

式~\eqref{eq:projection_decomp} の両辺に $\bnabla\cdot$ を作用させ，
非圧縮条件 $\bnabla\cdot\bu^{n+1}=0$ を代入すると $\bu^{n+1}$ が消去される：
\begin{equation}
  \bnabla^2q = \bnabla\cdot\bu^*
  \label{eq:why_divergence}
\end{equation}
この Poisson 方程式（PPE）を解いて $q$ を求め，$\bu^{n+1} = \bu^* - \bnabla q$ で速度を補正する．

\noindent\textbf{変密度への拡張：}
上記は等密度（$\rho = \mathrm{const}$）の場合である．
気液二相流では $\rho$ が界面をまたいで 2--3 桁変化するため，
速度補正は
$\bu^{n+1} = \bu^* - (\Delta t_{\mathrm{proj}}/\rho^{n+1})\bnabla\delta p$ となり，
PPE は可変係数の楕円型方程式
$\bnabla\cdot(1/\rho\,\bnabla\delta p)
= (1/\Delta t_{\mathrm{proj}})\bnabla\cdot\bu^*$ に変わる．
段階的な導出は §~\ref{sec:projection_derivation} で行う．

% paper/sections/08b_pressure.tex
% §8.2「変密度射影法の段階的導出」
% ═══════════════════════════════════════════════════════════════════════════════
\subsection{変密度射影法の段階的導出}
\label{sec:projection_derivation}
% ═══════════════════════════════════════════════════════════════════════════════

第~\ref{sec:collocate}章（§~\ref{sec:helmholtz}）では等密度場を仮定した Helmholtz 分解で
射影法の幾何学的意味を説明した．
本節ではこれを\textbf{変密度場}へ拡張し，
なぜ可変係数演算子 $\bnabla\cdot(1/\rho\,\bnabla p)$ が現れるかを 3 段階で導出する．

本稿では van Kan（1986）の \textbf{IPC（Incremental Pressure Correction）法}~\cite{vanKan1986}を採用し，
予測段に前時刻圧力 $-\bnabla p^n$ を組み込む
（§~\ref{sec:time_int}，式~\eqref{eq:predictor_imex_bdf2}）．
これにより標準 Chorin 法~\cite{Chorin1968}の $\Ord{\Delta t}$ 分割誤差が $\Ord{\Delta t^2}$ に低減される
（二次精度射影法の系統解説は~\cite{BellColella1989,Guermond2006}）．
以下では圧力射影部分だけを抽出し，射影に用いる時間幅を
$\Delta t_{\mathrm{proj}}$ と書く
（IMEX-BDF2 では $\Delta t_{\mathrm{proj}}=\gamma\Delta t$，
BDF1 startup では $\Delta t_{\mathrm{proj}}=\Delta t$；§\ref{sec:ipc_derivation}）．

\subsubsection{第 1 段階：予測速度}

射影段で圧力項だけを分離し，対流・粘性・外力の時間離散済み寄与を
$\bm{R}_{\mathrm{proj}}^{n+1}$ とまとめると：
\begin{equation}
  \rho^{n+1}\frac{\bu^{n+1} - \bu^n}{\Delta t_{\mathrm{proj}}}
  = \bm{R}_{\mathrm{proj}}^{n+1}
  - \bnabla p^{n+1}
  \label{eq:NS_full}
\end{equation}
IPC 法では前時刻圧力 $-\bnabla p^n$ を陽的に加えて\textbf{予測速度} $\bu^*$ を定義する：
\begin{equation}
  \rho^{n+1}\frac{\bu^* - \bu^n}{\Delta t_{\mathrm{proj}}}
  = \bm{R}_{\mathrm{proj}}^{n+1} - \bnabla p^n
  \quad\Longrightarrow\quad
  \bu^* = \bu^n
  + \frac{\Delta t_{\mathrm{proj}}}{\rho^{n+1}}
    \bigl(\bm{R}_{\mathrm{proj}}^{n+1} - \bnabla p^n\bigr)
  \label{eq:predictor}
\end{equation}
$\bu^*$ は一般に発散フリーではない（$\bnabla\cdot\bu^* \neq 0$）．

\subsubsection{第 2 段階：圧力補正}

式~\eqref{eq:NS_full} から式~\eqref{eq:predictor} を引くと，
圧力増分 $\delta p \equiv p^{n+1} - p^n$ と速度補正の関係が得られる：
\begin{equation}
  \rho^{n+1}\frac{\bu^{n+1} - \bu^*}{\Delta t_{\mathrm{proj}}} = -\bnabla(\delta p)
  \quad\Longrightarrow\quad
  \bu^{n+1} = \bu^* - \frac{\Delta t_{\mathrm{proj}}}{\rho^{n+1}}\bnabla(\delta p)
  \label{eq:corrector}
\end{equation}

\subsubsection{第 3 段階：発散フリー条件から PPE を導く}

$\bu^{n+1}$ に非圧縮条件 $\bnabla\cdot\bu^{n+1}=0$ を課す．
式~\eqref{eq:corrector} の両辺に発散を作用させると：
\begin{equation}
  \bnabla\cdot\!\left(\frac{1}{\rho^{n+1}}\bnabla(\delta p)\right)
  = \frac{1}{\Delta t_{\mathrm{proj}}}\,\bnabla\cdot\bu^*
  \label{eq:PPE_varrho}
\end{equation}
これが変密度 PPE（IPC 増分形式）である．
一様密度（$\rho=\mathrm{const}$）では
$\bnabla^2(\delta p) = \rho\,\bnabla\cdot\bu^*/\Delta t_{\mathrm{proj}}$ に帰着する．

\noindent\textbf{補足：}
外力 $\bm{F}_\text{ext}^{n+1}$ は重力項を含む．
表面張力の扱いは経路で分かれる：一括 smoothed-Heaviside PPE の縮約検証では
CSF 体積力 $(\kappa_{lg}^{n+1}\bnabla\psi^{n+1})/We$ を予測段に保持するが，
本稿の標準圧力ジャンプ経路では CSF 体積力としては入れず，
$j_{gl}=p_g-p_l=-\sigma\kappa_{lg}$ から作る毛管補正量を
第 3 段階の PPE/補正段に組み込む
（§~\ref{sec:split_ppe_pressure_jump_formulation}）．
密度 $\rho^{n+1}$ は，界面輸送により新時刻の配置として与えられている．
本節の導出は Backward Euler 形式で統一したが，
対流項は IMEX-BDF2（式~\eqref{eq:predictor_imex_bdf2}），
粘性項は implicit-BDF2（$\Ord{\Delta t^2}$，§~\ref{sec:time_viscous}）を用いる．
高密度比の分相 PPE 解法における圧力勾配の評価には
HFE による場延長を要する（§~\ref{sec:field_extension}，§~\ref{sec:split_ppe_recovery}）．

\noindent\textbf{PPE の離散化方針：}
以下の演算子形式は圧力増分 $\delta p$ および圧力 $p$ のいずれにも適用可能であるため，
一般性のため $p$ で表記する．
一括 PPE 縮約で式~\eqref{eq:PPE_varrho} の可変係数演算子を product-rule で展開すると：
\begin{equation}
  \bnabla\cdot\!\left(\frac{1}{\rho}\bnabla p\right)
  = \frac{1}{\rho}\bnabla^2 p - \frac{\bnabla\rho}{\rho^2}\cdot\bnabla p
  \label{eq:varrho_product_rule}
\end{equation}
この形式は高次 CCD 系の演算子で表せる（§~\ref{sec:ccd_poisson_matrix}）．
ただし高密度比・低 We の標準経路では，product-rule 展開した一括場ではなく，
分相圧力ジャンプ PPE と欠陥補正法（§~\ref{sec:defect_correction_ppe}）で
同じ有効面空間 $D_fG_A$ の固定点を定める．
有限体積法との等価性・調和平均の根拠は §~\ref{sec:fvm_face_coeff} 参照．
仮想時間刻みとゲージ固定の補足は付録~\ref{sec:ppe_pseudotime} 参照．

ただし，式~\eqref{eq:PPE_varrho} をどの線形系として解くかだけでは，
静止界面における力の釣り合いは保証されない．
PPE 内の圧力勾配，速度補正に用いる圧力勾配，および界面力が
同じ面評価位置と係数を共有することが，後続の Balanced--Force 条件の主題である．

\subsubsection{圧力反力としての面勾配：PPE だけでは代表が決まらない}
\label{sec:variational_pressure_reaction}

圧力射影を面速度空間で書くと，補正段は
\begin{equation}
  \bu_f^{n+1}
  =
  \bu_f^\ast
  -
  \Delta t_{\mathrm{proj}}\,
  a_f(p;c),
  \qquad
  a_f(p;c)=G_Ap-c_f
  \label{eq:subtractive_face_pressure_reaction}
\end{equation}
である．ここで $c_f$ は既知の毛管補正量または界面ジャンプ補正量であり，
$G_Ap$ は圧力が面速度へ及ぼす反力である．
この $G_A$ は，単に $D_fG_Ap$ が PPE 左辺を作るという条件だけでは決まらない．
圧力が非圧縮制約のラグランジュ乗数であるためには，面運動量計量 $M_A$ と
圧力節点計量 $W_p$ に対して
\begin{equation}
  \langle G_Ap,w_f\rangle_{M_A}
  +
  \langle p,D_fw_f\rangle_{W_p}
  =0
  \qquad
  \forall p,\ \forall w_f
  \label{eq:pressure_reaction_green_identity}
\end{equation}
を満たさなければならない（符号は式~\eqref{eq:subtractive_face_pressure_reaction}
の減算型補正に合わせたもの）．
したがって
\begin{equation}
  G_A=-M_A^{-1}D_f^TW_p,
  \qquad
  L_Ap\equiv D_fG_Ap
  \label{eq:variational_pressure_operator_pair}
\end{equation}
が，固定した $D_f,M_A,W_p$ に対する圧力反力とスカラー PPE の組である．
この組を用いると，圧力補正による仕事は
\begin{equation}
  \langle G_Ap,G_Ap\rangle_{M_A}
  =
  -\langle p,L_Ap\rangle_{W_p}
  \label{eq:pressure_reaction_energy_identity}
\end{equation}
として同じ離散内積で読める（$L_A$ の符号は
式~\eqref{eq:variational_pressure_operator_pair} の発散向きに従う）．

ここで重要なのは，$D_fG_Ap$ が同じでも $G_Ap$ の代表は一意ではない点である．
もし
\[
  \widetilde Gp = G_Ap+Zp,\qquad D_fZp=0,\qquad Zp\ne0
\]
なら，$\widetilde G$ は同じ PPE 右辺を満たし得るが，
式~\eqref{eq:pressure_reaction_green_identity} の圧力仕事を一般には満たさない．
この発散ゼロの余剰成分は，静的液滴では寄生流れとして，
非平衡界面では毛管駆動の過小または過大評価として現れる．
したがって本稿の圧力ジャンプ標準経路では，
補正段の圧力面加速度を
式~\eqref{eq:pressure_reaction_green_identity} の変分随伴代表に限定し，
高次コンパクト勾配のような発散同値な代表は，
この随伴性が証明された場合を除き，圧力反力としては用いない．
保存形共通フラックス経路では，この圧力射影の $D_f$ が
界面・密度・運動量輸送の $D_f$ でもなければならない．
射影が $D_f\bu_f=0$ を満たすという主張は，その同じ $D_f$ で
$q,m,\bm{p}_m$ が更新されるときにだけ質量輸送の非圧縮性を意味する．
したがって圧力反力の随伴性と共通フラックス台帳は，
どちらも同じ面速度空間を固定するための条件である．

\subsubsection{非一様格子での面平均補正}
\label{sec:fvm_ccd_corrector}

均一格子・周期境界条件では，PPE 内の勾配と速度補正の勾配に同じ
CCD 演算子を用いれば P-2（PPE 側 $G_h$ と補正側 $G_h$ の共有）が成立する．
一方，界面適合非一様格子と壁面境界条件では，PPE が面フラックス
$D_h(A_f G_h p)$ を有限体積法的に積分するのに対し，節点 CCD の速度補正は
節点メトリックで勾配を評価するため，面メトリックとの不一致が残る．

この経路では速度補正を面平均演算子
$\mathcal{G}^\text{adj}$ に切り替え，PPE の離散拡散項も
同じ随伴ペア $D_h^\text{bf}=-(\mathcal{G}^\text{adj})^\ast$ から組み立てる．
これにより PPE と補正が同じ面フラックス空間を共有し，
非一様格子での有限体積法--CCD メトリック不整合を避ける．
圧力ジャンプ条件を併用する場合も同じ面フラックス空間で
$G_{\Gamma}^{\mathrm{adj}}p=G^{\mathrm{adj}}p-B^{\mathrm{adj}}(j_{gl})$ とし，
$B^{\mathrm{adj}}(j_{gl})$ を面平均勾配と同じ局所物理距離 $H_f$ で定義する．
この面平均条件は本節で定義し，静止液滴での安定化効果は
§~\ref{sec:bf_static_droplet} で検証する．

\phantomsection
\label{sec:dccd_pressure_nofilt}% backward-compat: 外部参照保持


% ═══════════════════════════════════════════════════════════════════════════════
\subsection{チェッカーボード問題の解決策の概観}
\label{sec:checkerboard_solution}
% ═══════════════════════════════════════════════════════════════════════════════

コロケート格子の不安定性には\textbf{異なる 2 つの原因}が存在するため，
それぞれ独立した機構で対処する：
\textbf{面上の力の釣合不整合}（Balanced--Force 条件で解消）と
\textbf{格子スケールの圧力・速度分離}（同じ面空間による PPE/補正段共有で解消）である．
どちらか一方のみでは寄生流れまたはチェッカーボードが残存する．

コロケート格子における $2\Delta x$ チェッカーボード不安定の解消には，
歴史的には Rhie と Chow~\cite{RhieChow1983} による面速度補間が広く用いられてきたが，
本稿では採用しない．代わりに以下の二機構の組み合わせによりデカップリングを解消する：
\begin{enumerate}[noitemsep]
  \item \textbf{Balanced--Force 演算子整合}（§~\ref{sec:balanced_force}）：
        圧力勾配と表面張力を同じ面評価位置・同じ随伴ペアへ通すことで，
        演算子不整合に起因する寄生流れを構造的に除去
  \item \textbf{FCCD 面フラックス部分系}（§~\ref{sec:fccd_bf_sub_system}）：
        PPE 左辺，右辺発散，速度補正，毛管ジャンプ補正を
        同じ標準面空間上の $D_fG_A$ として組み立てる．
        標準圧力ジャンプ閉包では表面エネルギー仮想仕事に基づく毛管補正量と成分体積反力を
        同じ圧力仕事と随伴な面係数で構成し，
        純 FCCD DNS（§~\ref{sec:pure_fccd_dns}）では同じ面整合を全演算子へ拡張する
\end{enumerate}

% ═══════════════════════════════════════════════════════════════════════════════
\subsection{Balanced--Force 離散化：寄生流れ抑制のための離散化要件}
\label{sec:balanced_force}
% ═══════════════════════════════════════════════════════════════════════════════

\textbf{Balanced--Force 離散化条件}とは，圧力による面加速度
$A_f G_f p$ と毛管ジャンプによる既知の面上補正量 $c_f(j_{gl})$ を
\textbf{同一の離散評価位置と随伴整合した演算子対}で評価することで，
静止流体における Young--Laplace 平衡
$A_f G_f p=c_f(j_{gl})$ が離散化後にも同じ面空間で表現されることを
要求する条件である~\cite{Francois2006}．
CSF モデル（§~\ref{sec:csf}）を一括一流体経路として使う場合は，
この条件が $\bnabla p=(\kappa_{lg}/We)\bnabla\psi$ の同一評価位置での評価に還元される．
標準圧力ジャンプ経路では $j_{gl}=p_g-p_l=-\sigma\kappa_{lg}$ を
§~\ref{sec:split_ppe_capillary_range_projection} の面上補正量として扱う．
本節では，寄生流れ（第~\ref{sec:intro}章§~\ref{sec:failure_modes}）を
抑制するためのこの離散要件を定式化する．

% ─────────────────────────────────────────────────────────────────────────────
\subsubsection{寄生流れと計算爆発の根本原因：離散化演算子の不一致}
\label{sec:bf_operator_mismatch}
% ─────────────────────────────────────────────────────────────────────────────

静止液滴の平衡状態を考える．速度 $\bu=\bm{0}$ のとき，Navier--Stokes 方程式（式~\eqref{eq:NS_dimless_cls}）の運動量バランスは，圧力勾配と表面張力の完全な釣り合い：
\[
  \bnabla p = \frac{\kappa_{lg}\bnabla\psi}{We}
\]
を要求する．数値シミュレーションで寄生流れの数値成分を抑えるためには，
\textbf{離散化空間においてもこの等式を同じ演算子で近似}しなければならない．
離散化演算子を $G$ と表すと，
\[
  G_p(p) = G_\sigma\!\left(\frac{\kappa_{lg}\bnabla\psi}{We}\right)
\]
が要求される．

圧力勾配の評価 $G_p$ に高次精度 CCD 演算子 $D^{(1)}_{\mathrm{CCD}}$ を用い，
表面張力の空間微分 $G_\sigma$ に低次の有限差分 $D^{(1)}_{\mathrm{FD}}$（例：2次中心差分）を用いるような
非整合な離散化を考える．この場合，平衡条件 $\bnabla p = (\kappa_{lg}/We)\bnabla\psi$ のもとで
主要項が相殺した後の残差は：
\begin{equation}
  D^{(1)}_{\mathrm{CCD}} p - \frac{\kappa_{lg}}{We} D^{(1)}_{\mathrm{FD}}\psi
  = \Ord{h^2} \neq \bm{0}
  \label{eq:bf_operator_mismatch}
\end{equation}
となる（CCD の打ち切り誤差 $\Ord{h^6}$ と FD の $\Ord{h^2}$ の差 $\approx \Ord{h^2}$ が残差を支配する；付録~\ref{app:balanced_force_taylor} に導出）．この $\Ord{h^2}$ の不整合が静止しているべき流体を加速させる\textbf{人工的な体積力（駆動源）}として毎ステップ運動量方程式に現れる．固定外力であれば有界な応答しか生まないが，本問題では\textbf{正帰還ループ}が成立する：寄生速度が CLS 移流を通じてインターフェース形状を変形させ，変形した界面が曲率 $\kappa_{lg}$ の誤差を増幅し，さらに大きな表面張力の不整合を生む，という連鎖である．
特にコロケート格子では，変密度 PPE の離散演算子が $2h$ チェッカーボードモードに対してほぼ零固有値を持つ（§~\ref{sec:checkerboard_solution}）．
$\Ord{h^2}$ 残差の駆動力がこの零モード方向に投影されると，PPE 線形系の残差だけでは検出しにくいまま
補正速度に $O(1/\lambda_{\min})$ の増幅が生じる．
デカップリング対策が欠けていると，以下の正帰還ループが成立する：
\begin{itemize}
  \item 演算子不整合 $\Rightarrow$ 運動量方程式への人工体積力
  \item 人工体積力 $\Rightarrow$ 寄生速度の発生
  \item 寄生速度 $\Rightarrow$ CLS 移流による界面変形
  \item 界面変形 $\Rightarrow$ 曲率誤差増幅 $\Rightarrow$ さらに大きな演算子不整合
\end{itemize}
この増幅がステップごとに積算されると\textbf{指数的増大（計算爆発）}を招く．
一括 CSF 一流体経路では，Balanced--Force と FCCD 面フラックス部分系
（§~\ref{sec:fccd_bf_sub_system}）の組み合わせによりこのループを抑える
（§~\ref{sec:bf_static_droplet} 参照）．
一方，本稿の標準圧力ジャンプ経路では，表面張力を予測段階の CSF 体積力として入れず，
向き付き圧力ジャンプ $j_{gl}$ を切断面代表として使いながら，
表面エネルギーの仮想仕事から毛管補正量を構成し，
圧力仕事と随伴な成分体積制約付き鞍点系で静的反力と非平衡駆動を分離してから
PPE/補正段へ渡す（§~\ref{sec:split_ppe_capillary_range_projection}）．

本稿の定式化では，式~\eqref{eq:bf_operator_mismatch} の主要な演算子不整合を抑えるため，
圧力勾配・毛管補正量・速度補正を同じ面空間に置く．
一様格子の一括 CSF 経路では $D^{(1)}_\mathrm{CCD}$ 統一として読めるが，
非一様格子（$\alpha>1$）+ 壁面境界条件，および圧力ジャンプ標準経路では，
BF 部分系全体を面平均演算子 $\mathcal{G}^\text{adj}$ が定める
同じ面補正枠へ移し，PPE・速度補正・ジャンプ補正量を随伴ペアとして扱う
（詳細は §~\ref{sec:fvm_ccd_corrector}）．
この経路で節点 CCD の表面張力と $\mathcal{G}^\text{adj}$ の速度補正を混用すると
P-1/P-2 違反になるため採用しない．

\paragraph{2 経路構成の F-2 回避}
均一格子経路（$G_h = D^{(1)}_\mathrm{CCD}$ 統一）では，
PPE 内の $G_h p$ と速度補正の $G_h p$ が恒等的に同一演算子であるため
F-2（PPE 内 $G_h$ $\neq$ 補正 $G_h$）は構造的に回避される．
非一様格子経路では，BF 補正演算子対が $\mathcal{G}^\text{adj}$ に切り替わるため
\textbf{一見 F-2 違反に見える}が，PPE の離散拡散項にも同じ $\mathcal{G}^\text{adj}\circ\mathcal{G}^\text{adj}$
を adjoint として用いることで離散 BF 条件 $D_h = -G_h^\ast$（P-2）が保持される．
したがって両経路とも F-2 は回避され，差異は「どちらの演算子対を選ぶか」
（CCD ペア / 随伴ペア）に帰着する（面平均条件は §~\ref{sec:fvm_ccd_corrector}，静止液滴での検証は §~\ref{sec:bf_static_droplet}）．

\paragraph{Balanced--Force 条件（離散化要件のまとめ）}
CSF モデルを使用する場合，圧力勾配項 $-\bnabla p$ と表面張力項 $(\kappa_{lg}\bnabla\psi)/We$ に
\textbf{同一評価位置の離散演算子対}を適用する必要がある．
一様格子の一括 CSF 経路では $D_x^{(1)}, D_y^{(1)}$（第~\ref{sec:CCD}章）を統一使用する：
\begin{align}
  \text{圧力勾配：} & \quad
    -D_x^{(1)} p_{i,j}, \quad -D_y^{(1)} p_{i,j}
    \label{eq:balanced_pressure} \\
  \text{表面張力：} & \quad
    \frac{\kappa_{lg,i,j}}{We} D_x^{(1)} \psi_{i,j}, \quad
    \frac{\kappa_{lg,i,j}}{We} D_y^{(1)} \psi_{i,j}
    \label{eq:balanced_surface}
\end{align}
演算子の不一致は式~\eqref{eq:bf_operator_mismatch} の $\Ord{h^2}$ 残差を生じ，
正帰還ループを通じて計算の発散を招く．
\textbf{本稿の標準圧力ジャンプ定式化では，
$A_f G_f p$ と表面エネルギー仮想仕事に基づく毛管補正量を同じ面空間で評価し，
成分補正済みの補正量を速度補正へ渡す
（§~\ref{sec:split_ppe_capillary_range_projection}）．}
一括 CSF 経路は，この面釣合原理を
$\bnabla p$ と $(\kappa_{lg}/We)\bnabla\psi$ の同一評価位置での評価として読んだ場合に対応する．

% paper/sections/08_0_bf_failure.tex
%
% §7.2.0  BF 失敗モード 5 種
%         §7.2（Balanced--Force）冒頭に配置し、
%         以降の P-1..P-7 原則（§7.3）の動機付けとする。
%
% ═══════════════════════════════════════════════════════════════════════════════
\subsection{Balanced--Force 失敗モード 5 種}
\label{sec:bf_failure_modes}
% ═══════════════════════════════════════════════════════════════════════════════

連続問題での静止液滴平衡
\begin{equation}
  -\bnabla p + \mathbf{f}_\sigma = \mathbf{0}
  \quad\text{（Young--Laplace: }j_{gl}=p_g-p_l=-\sigma\kappa_{lg}\text{）}
  \label{eq:bf_continuous_balance}
\end{equation}
を\textbf{離散的に}保持するには
\begin{equation}
  -G_h p + f_{\sigma,h} \approx \mathbf{0}
  \label{eq:bf_discrete_balance}
\end{equation}
が要請される．微分階数ではなく\textbf{離散演算子の整合性}が
寄生流れ（spurious currents）を支配する~\cite{Popinet2009,Renardy2002,DennerVanWachem2015}．
離散化で現れる 5 つの破綻モードを以下に列挙する：

\begin{tcolorbox}[warnbox,title={F-1..F-5：Balanced--Force 失敗モード}]
\begin{description}[leftmargin=2em,style=nextline]
  \item[\textbf{F-1} 評価位置不一致]
    $G_h p$ と $f_{\sigma,h}$ が異なる格子位置（節点 vs 面）で評価される．
    $(G_h p)_f \neq (f_\sigma)_f$ が定義上保証できない．
  \item[\textbf{F-2} PPE 勾配と補正勾配の不一致]
    PPE 内部で用いる $G_h$ と，$\bm{u}^* \to \bm{u}^{n+1}$ 補正で用いる $G_h$ が異なる．
    ある方程式を解き，別の方程式で補正していることになる．
  \item[\textbf{F-3} $D_h$ と $G_h$ が離散共役でない]
    $D_h \neq -G_h^\ast$．PPE 演算子が SPD でなくなり，
    離散エネルギー恒等式を崩し，CG が適用できない．
  \item[\textbf{F-4} $\betaf=(1/\rho)_f$ の経路不一致]
    PPE / 速度補正 / 界面張力の 3 経路で $\betaf$ の平均化
    （算術 vs 調和）が異なる．高密度比 $\rho_l/\rho_g\gg 1$ で
    補正残差が界面上で $\Ord{1}$ となる．
  \item[\textbf{F-5} 界面ジャンプ無補正の CCD ステンシル]
    $j_{gl}=-\sigma\kappa_{lg}$ を跨ぐ CCD ステンシルが
    滑らかな場として扱われ，ジャンプが数値的に平滑化される．
    GFM/IIM による面フラックス補正が必要．
\end{description}
\end{tcolorbox}

\noindent\textbf{重要な事実}：
F-1..F-5 は\textbf{すべて微分階数と独立}である．
「高次 CCD を全量に適用」した素朴離散化は F-1..F-5 のどれか（多くの場合複数）を
必ず踏む．本稿の rising bubble PoC では F-2 と F-4 が同時発生し，
PPE RHS の発散だけを修正しても補正側が別 $G_h$ を参照する限り
BF 残差は減らない（P-2 参照，\S\ref{sec:bf_seven_principles}）．

\noindent\textbf{方針}：
§\ref{sec:bf_seven_principles} で 7 つの設計原則 P-1..P-7 を定式化し，
§\ref{sec:fccd_bf_sub_system} で FCCD 面ジェットによる
面フットプリント不一致の $\Ord{H^4}$ 抑制と，
GFM/CSF 律速を含む実効 BF 残差の $\Ord{H^2}$ 評価を分離して示す．

% ═══════════════════════════════════════════════════════════════════════════════

% paper/sections/08_1_bf_seven_principles.tex
%
% §7.3  Balanced-Force 七原則
%       P-1 位置一致 / P-2 PPE-corrector 共有 / P-3 geometry 経路統一 /
%       P-4 β_f 統一 / P-5 曲率への CCD 非適用 /
%       P-6 界面ジャンプ補正 / P-7 BF sub-system 分離
%
% ═══════════════════════════════════════════════════════════════════════════════
\subsection{Balanced-Force 七原則}
\label{sec:bf_seven_principles}
% ═══════════════════════════════════════════════════════════════════════════════

§\ref{sec:bf_failure_modes} の 5 失敗モードを回避するため，
本節では 7 原則 P-1..P-7 を与える．
\textbf{BF 整合性の核は P-1（位置一致）/ P-2（補正 $G_h$ 共有）/ P-4（$\betaf$ 統一）/
P-6（界面ジャンプ補正）の 4 原則}が担い，これらが欠けると寄生流れが
$\Ord{h^0}$ で残留する．残る P-3（geometry 経路統一）/ P-5（曲率非 CCD 化）/
P-7（sub-system 分離）は\textbf{次数改善および予防的役割}を持つ．
\textbf{特に P-5 の status は適用形態に依存}：対称フィルタ適用時は
P-1 維持の必要条件として連動する（§\ref{sec:bf_p5_curvature}），
非対称・無フィルタ時は次数改善の予防策に留まる．主解消は依然 4 原則 (P-1/P-2/P-4/P-6) で完結する
（詳細は §\ref{sec:bf_p5_curvature} および §\ref{sec:bf_failure_principle_map} の F$\times$P 表参照）．

% -------------------------------------------------------------------------------
\subsubsection{\texorpdfstring{P-1：$G_h p$ と $f_{\sigma,h}$ の同一位置評価}{P-1：Gh p と f sigma h の同一位置評価}}
\label{sec:bf_p1_location}
% -------------------------------------------------------------------------------

\textbf{コロケート配置の面演算子 FCCD（\S\ref{sec:fccd_def}）で実装する場合}，
速度補正は面値 $u_f^{n+1}$ を経由して節点に戻す 2 段設計となる：
\begin{equation}
  u_f^{n+1} = u_f^* + \Delta t_{\mathrm{proj}}\Bigl(
    -\betaf(G_h^\text{bf} p)_f + (f_\sigma)_f + \cdots
  \Bigr),
  \label{eq:bf_p1_corrector}
\end{equation}
$(G_h^\text{bf} p)_f$ と $(f_\sigma)_f$ は\textbf{同じ面で評価}されねばならない．
節点 CCD 勾配を面に補間し，$f_\sigma$ を別の補間で構築する
\textbf{「位置を合わせれば十分」は不正解}．
補間演算子自身が互いに随伴整合（adjoint-compatible）でなければ
$(G_h p)_f - (f_\sigma)_f$ は相殺しない．
P-1 から P-7 まで本節全体を通じ，「面」は FCCD で定義される共通面ジェット
位置（\S\ref{sec:face_jet_def}）を指す．

% -------------------------------------------------------------------------------
\subsubsection{\texorpdfstring{P-2：随伴 PPE による F-2・F-3 の同時解消：$D_h^\text{bf} = -(G_h^\text{bf})^\ast$}{P-2：随伴 PPE による F-2・F-3 の同時解消}}
\label{sec:bf_p2_adjoint}
% -------------------------------------------------------------------------------

PPE を\textbf{随伴ペア}から組み立てる：
\begin{equation}
  D_h^\text{bf}\Bigl(\betaf\,G_h^\text{bf}\,p\Bigr)
    = \frac{1}{\Delta t_{\mathrm{proj}}}\,D_h^\text{bf}\,\mathbf{u}^*,
  \qquad D_h^\text{bf} = -(G_h^\text{bf})^\ast.
  \label{eq:bf_p2_ppe}
\end{equation}
これにより離散内積 $\langle\cdot,\cdot\rangle_h$ 下で
\begin{equation}
  \langle p,\,A_h p\rangle_h
    = \langle G_h^\text{bf} p,\,\betaf\,G_h^\text{bf} p\rangle_h
    \ge 0,
  \label{eq:bf_p2_spd}
\end{equation}
零空間（定数モード）除去後に $A_h$ は SPD で CG が主解法，
\textbf{補正側 $G_h$ = PPE 側 $G_h^\text{bf}$} が保証される．

\noindent\textbf{rising bubble への含意}：
PPE RHS の発散だけを修正しても，補正側が別の $G_h$ を参照する限り
F-2 は解消しない．$G_h^\text{bf}$ を両経路で共有することが必要．

\noindent\textbf{F-2 と F-3 の独立性}：
F-3（$D_h\neq -G_h^\ast$；随伴ペア違反）と
F-2（PPE 内 $G_h\neq$ 補正 $G_h^\text{bf}$；経路違反）は\textbf{独立の失敗モード}だが，
P-2 の随伴構築 $D_h^\text{bf}=-(G_h^\text{bf})^\ast$ は両者を同時に解消する：
$G_h^\text{bf}$ を固定して $D_h^\text{bf}$ をその随伴として導出するため，
PPE と補正の両経路で同一演算子が現れ，かつ随伴整合も自動的に成立する．

% -------------------------------------------------------------------------------
\subsubsection{\texorpdfstring{P-3：$f_\sigma$ を $[p]_\Gamma$ と同じ離散幾何に通す}{P-3：f sigma を圧力ジャンプと同じ離散幾何に通す}}
\label{sec:bf_p3_geometry}
% -------------------------------------------------------------------------------

CSF 形式の $f_\sigma$ を，仮想毛管圧 $p_\sigma$ から
$G_h^\text{bf} p_\sigma$ として構成できるよう設計すれば，
静止平衡は
\begin{equation}
  -G_h^\text{bf} p + G_h^\text{bf} p_\sigma
    = G_h^\text{bf}(p_\sigma - p) \approx \mathbf{0}
  \label{eq:bf_p3_cancel}
\end{equation}
が\textbf{PPE 解の精度}により成立する．
独立に離散化された 2 ベクトルのキャンセルを期待する旧来設計と本質的に異なる．

% -------------------------------------------------------------------------------
\subsubsection{\texorpdfstring{P-4：$\betaf$ を 3 経路で統一}{P-4：beta f を 3 経路で統一}}
\label{sec:bf_p4_beta}
% -------------------------------------------------------------------------------

変密度 PPE $L_h p = -D_h^\text{bf}(\betaf\,G_h^\text{bf} p)$ において
$\betaf = (1/\rho)_f$ は以下の 3 経路\textbf{すべてで同一}：
\begin{enumerate}[leftmargin=2em,noitemsep]
  \item PPE 演算子 $L_h$ 自身，
  \item 速度補正 $u_f^{n+1} = u_f^* - \Delta t_{\mathrm{proj}}\,\betaf\,(G_h^\text{bf} p)_f$，
  \item 界面張力正規化 $f_\sigma \propto \betaf\,\sigma\kappa\,(\nabla\psi)_f$（CSF
        標準式 $\sigma\kappa\nabla\psi$ に $\betaf$ 因子を付加した\textbf{変密度
        corrector 離散化の特殊形}；等密度では $\betaf$ が定数で標準式に還元する）．
\end{enumerate}
PPE に調和平均，補正に算術平均という「混用」は
$\rho_l/\rho_g \gg 1$ で BF 残差を界面で $\Ord{1}$ とする．
\textbf{本稿の選択}：$\betaf^\text{harm}=2/(\rho_L+\rho_R)$ を 3 経路共通で採用．

\smallskip\noindent\textbf{命名と等価性（混乱回避）}：
$a\equiv 1/\rho$ と置くと，FVM フェイス係数は
\textbf{$a$ の調和平均}
\begin{equation}
  \betaf^\text{harm}
  = a_f^\text{harm}
  = \frac{2a_La_R}{a_L+a_R}
  = \frac{2}{\rho_L+\rho_R}
  \label{eq:bf_beta_harmonic}
\end{equation}
である（§\ref{app:fvm_face_coeff}）．
これは 2 値界面では\textbf{密度の算術平均の逆数}
$[(\rho_L+\rho_R)/2]^{-1}$ と同じ値を与える．
一方，\textbf{$\rho$ の調和平均の逆数}
$(\rho_L+\rho_R)/(2\rho_L\rho_R)$ は別物であり，
高密度比では界面係数を誤る典型的な実装ミスである．

% -------------------------------------------------------------------------------
\subsubsection{P-5：曲率に生 CCD を適用しない}
\label{sec:bf_p5_curvature}
% -------------------------------------------------------------------------------

$\kappa = \bnabla\cdot(\bnabla\psi/|\bnabla\psi|)$ は
界面から $\Ord{\varepsilon}$ 幅でしか滑らかでない場を 2 回微分する．
高次 CCD は分母の $|\bnabla\psi|$ 小のノイズを増幅し，
ほぼ定数の bulk を広いステンシルで走査する．

\noindent\textbf{規則}：$\bnabla\psi$ には CCD を用いてよい．
$\kappa$ は\textbf{界面帯に制限した平滑化}または
2 次 FD に落とす．曲率精度は\textbf{演算子位置整合に比べ 2 次的}であり，
BF 残差への影響は 2 次効果に留まる．

% -------------------------------------------------------------------------------
\subsubsection{P-6：界面越えステンシルへのジャンプ補正}
\label{sec:bf_p6_jump}
% -------------------------------------------------------------------------------

CCD は滑らかな場を仮定する．静止液滴で $\Delta p = \sigma\kappa$ が
正しく設定されているのに寄生流れが出る場合，
ほぼ常に\textbf{界面越えステンシルが圧力ジャンプを平滑化}している．
必要な補正：
\begin{itemize}[noitemsep]
  \item \textbf{GFM 面フラックス補正}（Phase 1 実用）：
  \begin{equation}
    (G_h^\text{bf} p)_f^\text{corrected}
      = (G_h^\text{bf} p)_f - \frac{\sigma\kappa}{H_f},
    \label{eq:bf_p6_gfm}
  \end{equation}
  bulk ステンシルは変更せず，ジャンプだけを面フラックスに注入．
  2 次精度；Phase 1 検証と H-01 診断に十分．
  \item \textbf{IIM コンパクト行補正}（Phase 3 研究）：
  FCCD のコンパクト関係 $M_f\mathbf{g} = R_f\mathbf{p}$ に対し
  $[p]_\Gamma,\,[\betaf\partial_n p]_\Gamma$ を Taylor 展開した
  補正ベクトル $c_\Gamma^\text{IIM}$ を加算．
\end{itemize}

% -------------------------------------------------------------------------------
\subsubsection{P-7：BF sub-system と bulk 微分子の分離}
\label{sec:bf_p7_separation}
% -------------------------------------------------------------------------------

\begin{center}
\begin{tabular}{@{}lll@{}}
\hline
タスク & 演算子 & 精度要求 \\
\hline
bulk 移流・拡散（$u,v$）& 節点 CCD / UCCD6 / FCCD & 高次（$\Ord{h^6}$） \\
PPE + 補正 & $G_h^\text{bf},\,D_h^\text{bf}$（面随伴ペア）& 内整合（$\Ord{h^2}$--$\Ord{h^4}$） \\
界面張力 $f_\sigma$ & $G_h^\text{bf}$ と同じ面スロット & 内整合 \\
曲率 $\kappa$ & 界面帯制限平滑化 & 2 次 \\
\hline
\end{tabular}
\end{center}

\noindent\textbf{設計哲学}：
BF sub-system は\textbf{高次である必要はない}．
内部整合性（位置一致・随伴ペア・$\betaf$ 統一）があれば
寄生流れは演算子精度ではなく
\textbf{界面ジャンプ処理の精度}で律速される．
bulk と BF で演算子を分けるのは「妥協」ではなく\textbf{設計原則}．

% -------------------------------------------------------------------------------
\subsubsection{F-1--F-5 × P-1--P-7 対応マップ}
\label{sec:bf_failure_principle_map}
% -------------------------------------------------------------------------------

5 失敗モード（§\ref{sec:bf_failure_modes}）と 7 原則の対応関係：
\begin{center}
\renewcommand{\arraystretch}{1.2}
\begin{tabular}{@{}lcccccccl@{}}
\hline
 & P-1 & P-2 & P-3 & P-4 & P-5 & P-6 & P-7 & 主解消原則 \\
\hline
F-1（位置ズレ）              & $\blacksquare$ & & $\square$ & & & & $\square$ & P-1 \\
F-2（補正 $G_h$ 不一致）      & & $\blacksquare$ & & & & & $\square$ & P-2 \\
F-3（随伴違反）              & & $\blacksquare$ & & & & & & P-2 \\
F-4（$\betaf$ 混用）         & & $\square$ & & $\blacksquare$ & & & & P-4 \\
F-5（界面越えステンシル）\footnote{F-5 で P-6 単独は \(\Ord{h^2}\) を達成する；P-3 と P-5 は次数改善の上乗せ寄与であり主解消ではない．}   & & & $\square$ & & $\square$ & $\blacksquare$ & & P-6 \\
\hline
\end{tabular}

{\footnotesize $\blacksquare$：主解消，$\square$：副次的に寄与}
\end{center}

\noindent P-5（曲率 $\kappa$ の非 CCD 化）と P-3（$f_\sigma$ を $p$ と同幾何）は
F-5（界面越えステンシル平滑化）に対し副次的・間接的に寄与する．
P-7（sub-system 分離）は F-1・F-2 の\textbf{設計レベル}での発生を防ぐ予防策．
\textbf{F-1..F-5 への主解消は P-1, P-2, P-4, P-6 の 4 原則}が担い，
P-3, P-5, P-7 は次数改善・設計予防として\textbf{独立な設計目的}を持つ
（各原則の設計動機は §\ref{sec:bf_p1_location}--\ref{sec:bf_p7_separation} 参照）．
7 原則の動機は単一の「5 モード解消最小集合」ではなく，
主解消 4 原則と副次寄与 3 原則が\textbf{独立な目的群}を網羅する設計である．

% ═══════════════════════════════════════════════════════════════════════════════

% paper/sections/08_2_fccd_bf.tex
%
% §8.5  FCCD 面フラックス部分系の構成
%
% ═══════════════════════════════════════════════════════════════════════════════
\subsection{FCCD 面フラックス部分系の構成}
\label{sec:fccd_bf_sub_system}
% ═══════════════════════════════════════════════════════════════════════════════

§~\ref{sec:fccd_def} で導出した FCCD 面フラックス表現
$\FaceJet{q}=(q_f,q'_f,q''_f)$ は，BF 部分系の
P-1..P-4 を\textbf{同時に}満たす面表現である．
圧力勾配，速度補正，HFE/GFM ジャンプ行，毛管 cochain を同じ面フットプリントへ写すことで，
連続平衡
$-\bnabla p+\mathbf{f}_\sigma=\mathbf{0}$
を離散空間でも同じ面釣合として扱える．

% -------------------------------------------------------------------------------
\phantomsection
\label{sec:fccd_face_residual_order}

本稿の構成では {\FCCD} 面フラックス表現 $\FaceJet{q}=(q_f,q'_f,q''_f)$ が
圧力勾配・速度補正・HFE 片側データ・圧力ジャンプ条件の 4 経路で共有される（(R1)-(R4)）．
これにより面フットプリントは 4 経路で一致し，
\begin{equation}
  \mathcal{G}p := D^\FCCD p,
  \qquad c_f(j_{gl}) := A_f B_f(j_{gl})
  \qquad\Longrightarrow\qquad
  a_f(p;c)=A_f\,\mathcal{G}p-c_f
  \label{eq:bf_fccd_unification}
\end{equation}
となる．Balanced--Force の判定対象は，個別の節点勾配ではなく
同じ face complex 上の実効補正 cochain $a_f$ である．
標準 pressure-jump 閉包では $c_f$ を
§~\ref{sec:split_ppe_capillary_range_projection} の
surface-energy virtual-work/Riesz 構成で定め，
同じ面複体の Hodge gate で静的平衡か非平衡駆動かを判定する．
圧力勾配値域への射影はこの gate の診断代表であり，
本番補正前に divergence-free 毛管 cochain を一律に消すための操作ではない．

% -------------------------------------------------------------------------------
\subsubsection{\texorpdfstring{面フラックス表現 $\FaceJet{p}$ による BF 構築}{面フラックス表現 FaceJet(p) による BF 構築}}
\label{sec:face_jet_bf_construction}
% -------------------------------------------------------------------------------

面フラックス表現 $\FaceJet{p}=(p_f,p'_f,p''_f)$ は
P-1..P-4 を同一の面表現で具現する．
変密度 PPE の\textbf{面フラックス}は
\begin{equation}
  F^p_f := A_f\,\FaceJet{p}_1 = A_f\,p'_f,
  \label{eq:fccd_bf_face_flux}
\end{equation}
ここで $A_f$ は逆密度，切断面係数，局所距離，メトリクスを含む
P-4 の単一面係数である（一様格子では $\betaf=(1/\rho)_f$ に縮約）．
PPE は FVM セル積分で
\begin{equation}
  F^p_{i+1/2} - F^p_{i-1/2}
    = \int_{x_{i-1/2}}^{x_{i+1/2}}\! S\,dx,
  \qquad S = \frac{1}{\Delta t_{\mathrm{proj}}}\,D_h^\text{bf}\,\bu^*,
  \label{eq:fccd_bf_cell_balance}
\end{equation}
として組み立て，速度補正は
\begin{equation}
  u_f^{n+1}
    = u_f^* - \Delta t_{\mathrm{proj}}\,A_f\,p'_f
    \quad\text{（同じ }p'_f\text{）．}
  \label{eq:fccd_bf_corrector}
\end{equation}
\textbf{同じ $p'_f$} が PPE 組立と補正で共有されるため P-2 が構造的に成立．

% -------------------------------------------------------------------------------
\subsubsection{界面ジャンプ GFM 補正（標準 pressure-jump 閉包）}
\label{sec:fccd_bf_gfm_integration}
% -------------------------------------------------------------------------------

界面 $\Gamma$ を跨ぐ面 $f$ で面フラックス表現を
ジャンプ条件 $j_{gl}=p_g-p_l=-\sigma\kappa_{lg}$，$[A_f\partial_n p]_\Gamma=0$
（非相変化）で補正する．標準経路では，ジャンプ演算子
$B_f(j_{gl})$ を同じ局所物理距離 $H_f$ と同じ面係数 $A_f$ から定義し，
\begin{equation}
  a_f(p;c)
    = A_f G_f p - A_f B_f(j_{gl})
  \label{eq:fccd_bf_gfm_face}
\end{equation}
を補正段階の実効 face cochain とする．
相内の滑らかな離散化は保ち，ジャンプ項だけを面フラックスへ加える．
HFE はこの GFM/affine-jump 行に必要な片側 Hermite データを供給し，
$c_f=A_f B_f(j_{gl})$ または表面エネルギーから得る
Riesz cochain は同じ face slot で速度補正へ渡す．
本稿の標準構成では GFM/affine-jump + HFE + DC 残差契約 +
変分毛管 Hodge gate を用いる．
固定回数 $k=3$ は §~\ref{sec:defect_correction_ppe} の成分検証での実用上限であり，
補正段階の標準契約は高次残差が規定許容値へ落ちることである．

% -------------------------------------------------------------------------------
\subsubsection{静水圧テストによる BF 精度検証}
\label{sec:fccd_bf_hydrostatic}
% -------------------------------------------------------------------------------

\noindent\textbf{検証問題}：重力 $\mathbf{g}=-g\hat{\mathbf{z}}$ 下の静止 2 層流体．
解析解は $p_\text{stat}(z) = -\rho(\phi(z))\,g\,z$，$\bu=\bm{0}$．

\noindent\textbf{BF 残差測定}：
重力加速度を $\bm{a}^g_f$，静水圧の面勾配を
$(G_h^\text{bf}p_\text{stat})_f$ として，
一ステップ補正で速度へ注入される離散加速度残差を測る：
\begin{equation}
  \varepsilon_\text{BF}
    := \Delta t_{\mathrm{proj}}
       \max_f\,
       \Bigl|\,
         \bm{a}^g_f
         - A_f\,(G_h^\text{bf}p_\text{stat})_f
       \Bigr|.
  \label{eq:fccd_bf_hydrostatic_metric}
\end{equation}
これは $\bu^*=\bm{0}$ から開始し，重力加速度と
圧力補正を同じ面スロットで同時に適用した
1 ステップ釣合確認後の $\max_f|u_f^{n+1}|$ と同値である
（§~\ref{sec:time_buoyancy} の
balanced-buoyancy 分解と同じ考え方）．
理想離散 BF では $\varepsilon_\text{BF}=0$．このテストの読み方は
数値値そのものではなく，P-1/P-2/P-4 のいずれを破ったかを
一ステップ速度残差として分離できる点にある：
\begin{center}
\begin{tabular}{@{}p{0.30\textwidth}p{0.45\textwidth}p{0.18\textwidth}@{}}
\hline
構成 & 検出する違反 & 期待される判定 \\
\hline
非整合な節点勾配 + 事後面平均 & 面ロカス不一致，補正 $G_h$ 不一致 & FAIL \\
同一 face-adjoint pair & PPE と補正の面複体共有 & PASS \\
pressure-jump Hodge gate & 静的平衡でない毛管 Hodge 成分の消去 & FAIL \\
\hline
\end{tabular}
\end{center}
\noindent pressure-jump 経路の実測値は §~\ref{sec:bf_static_droplet} の
face-balance gate で示す．そこでは raw 毛管 cochain の非勾配成分と，
静的平衡で消えるべき成分・非平衡界面の物理駆動成分を分けて評価する．

\noindent\textbf{重力 BF 検証}：
本表は重力平衡での BF 残差を測定する
（CSF BF の静止液滴テスト §~\ref{sec:bf_static_droplet} とは独立；
CSF 側は $f_\sigma$ と $G_h p_\sigma$ の相殺が主論点）．
この診断の検証は §~\ref{sec:verify_split_ppe_dc} と §~\ref{sec:bf_static_droplet} に分離する．

詳細検証は §~\ref{sec:verify_split_ppe_dc}（要素検証）と
§~\ref{sec:bf_static_droplet}（静止液滴 gate）．

% -------------------------------------------------------------------------------
\subsubsection{面演算子閉包 A/B/C の位置付け}
\label{sec:fccd_bf_strategy_table}
% -------------------------------------------------------------------------------

\begin{center}
\begin{tabular}{@{}p{0.20\textwidth}p{0.40\textwidth}p{0.30\textwidth}@{}}
\hline
構成 & 内容 & 用途 \\
\hline
\textbf{A：標準 affine-jump face complex} &
$D_fG_A$，$B_f(j_{gl})$，変分毛管 Riesz cochain，
および component saddle / Hodge gate を同じ pressure-adjoint face complex で構築．
HFE は片側 Hermite データを供給．
& 高密度比・低 We の標準 pressure-jump 閉包．
\\
\textbf{B：CSF/BF 縮約} &
CSF 体積力を予測段へ入れ，PPE + 補正 + $f_\sigma$ を同じ面層に置く．
& 低密度比・一括 PPE の縮約検証経路．
\\
\textbf{C：FCCD DNS 極限} &
BF 面整合を CLS，運動量，分相 PPE，HFE/GFM へ広げる．
& 標準経路の高次極限としての研究構成．
\\
\hline
\end{tabular}
\end{center}

\noindent\textbf{本稿の標準構成}：
バルクの高次 CCD 系演算子で移流・拡散精度を確保しつつ，
PPE，速度補正，HFE/GFM ジャンプ行，毛管 face cochain は同じ BF 面演算子へ集約する．
標準 pressure-jump 閉包では affine-jump face law と変分毛管 Hodge gate を用い，
CSF/BF は低密度比・一括 PPE の縮約診断として扱う．
純 FCCD DNS はこの面整合を全演算子へ広げた高次極限として位置付ける
（§~\ref{sec:pure_fccd_dns}）．

% -------------------------------------------------------------------------------
\subsubsection{BF アンチパターン一覧}
\label{sec:fccd_bf_anti_patterns}
% -------------------------------------------------------------------------------

\begin{tcolorbox}[warnbox,title={避けるべき BF アンチパターン}]
\begin{itemize}[leftmargin=1.5em,noitemsep]
  \item 高次 $G_h p$，低次毛管 cochain（切断誤差非対称）
  \item 節点 CCD 全量 $\kappa_{lg}$（界面帯ノイズ増幅）
  \item PPE $G_h$ と補正 $G_h$ が別演算子（F-2）
  \item 毛管 cochain が節点，$G_h p$ が面（F-1）
  \item $A_f/\betaf$ 混用（PPE 調和 / 補正算術 / ジャンプ cochain 別メトリクス；F-4）
  \item 界面越え CCD ステンシル無補正（F-5）
  \item 表面エネルギーの仮想仕事を経ない raw 曲率 cochain を物理力として採用する
  \item 非平衡界面の毛管 Hodge 成分を range projection で本番補正前に消す
\end{itemize}
\end{tcolorbox}

% ═══════════════════════════════════════════════════════════════════════════════


\medskip
以上で本章の定式化を完了した．
コロケート格子・射影法の幾何学的基礎（§~\ref{sec:helmholtz}），
変密度への拡張（§~\ref{sec:projection_derivation}），
Balanced--Force 条件（§~\ref{sec:balanced_force}），
失敗モードと七原則（§~\ref{sec:bf_failure_modes}--\ref{sec:bf_seven_principles}），
FCCD 面フラックス部分系（§~\ref{sec:fccd_bf_sub_system}）がそろって初めて
§~\ref{sec:ppe_solve}章（次章）の PPE 求解が，力の釣り合いを保つ射影として意味を持つ．
次章では，§~\ref{sec:projection_derivation} で導出した可変係数 PPE を
いかに高精度・低コストで解くかを論じる．
